\documentclass{article}
\usepackage[utf8]{inputenc}
\usepackage[spanish]{babel}
\usepackage{amsmath}
\usepackage{amsfonts}
\usepackage{geometry}
\geometry{a4paper, margin=2.5cm}

\title{Soluciones: Distancia entre una recta y un punto (Ejercicios 29 al 35)}
\author{Fabio Hernan Realpe Martinez}
\date{4 de Agosto de 2026}

\begin{document}

\maketitle

\section*{Fórmula General}
La distancia $d$ desde un punto $P(x_1, y_1)$ a una recta $L$ dada por la ecuación general $Ax + By + C = 0$ se define como:
\begin{equation}
    d = \frac{|Ax_1 + By_1 + C|}{\sqrt{A^2 + B^2}}
\end{equation}

\section*{Soluciones a los Ejercicios}

\textbf{29.} Encontrar la distancia entre $2x - 3y = 4$ y el punto $(-7, -2)$.\\
Reescribimos la recta a la forma general: $2x - 3y - 4 = 0$. Aquí $A=2, B=-3, C=-4$.
\[
d = \frac{|2(-7) - 3(-2) - 4|}{\sqrt{(2)^2 + (-3)^2}} = \frac{|-14 + 6 - 4|}{\sqrt{4 + 9}} = \frac{|-12|}{\sqrt{13}} = \frac{12}{\sqrt{13}} = \frac{12\sqrt{13}}{13}
\]

\vspace{0.5cm}
\textbf{30.} Encontrar la distancia entre $-5x + 6y = 2$ y el punto $(1, 3)$.\\
Reescribimos la recta: $-5x + 6y - 2 = 0$. Aquí $A=-5, B=6, C=-2$.
\[
d = \frac{|-5(1) + 6(3) - 2|}{\sqrt{(-5)^2 + (6)^2}} = \frac{|-5 + 18 - 2|}{\sqrt{25 + 36}} = \frac{|11|}{\sqrt{61}} = \frac{11}{\sqrt{61}} = \frac{11\sqrt{61}}{61}
\]

\vspace{0.5cm}
\textbf{31.} Encontrar la distancia entre $2x - 4y = -42$ y el punto $(7, -21)$.\\
Reescribimos la recta: $2x - 4y + 42 = 0$. Aquí $A=2, B=-4, C=42$.
\[
d = \frac{|2(7) - 4(-21) + 42|}{\sqrt{(2)^2 + (-4)^2}} = \frac{|14 + 84 + 42|}{\sqrt{4 + 16}} = \frac{|140|}{\sqrt{20}} = \frac{140}{2\sqrt{5}} = \frac{70}{\sqrt{5}} = 14\sqrt{5}
\]

\vspace{0.5cm}
\textbf{32.} Encontrar la distancia entre $7x + 5y = 6$ y el punto $(0, 0)$.\\
Reescribimos la recta: $7x + 5y - 6 = 0$. Aquí $A=7, B=5, C=-6$.
\[
d = \frac{|7(0) + 5(0) - 6|}{\sqrt{(7)^2 + (5)^2}} = \frac{|-6|}{\sqrt{49 + 25}} = \frac{6}{\sqrt{74}} = \frac{6\sqrt{74}}{74} = \frac{3\sqrt{74}}{37}
\]

\vspace{0.5cm}
\textbf{33.} Encontrar la distancia entre $3x + 7y = 0$ y el punto $(-2, -8)$.\\
La recta ya está en forma general: $3x + 7y + 0 = 0$. Aquí $A=3, B=7, C=0$.
\[
d = \frac{|3(-2) + 7(-8)|}{\sqrt{(3)^2 + (7)^2}} = \frac{|-6 - 56|}{\sqrt{9 + 49}} = \frac{|-62|}{\sqrt{58}} = \frac{62}{\sqrt{58}} = \frac{31\sqrt{58}}{29}
\]

\vspace{0.5cm}
\textbf{34.} Encontrar la distancia entre $11x - 12y = 5$ y el punto $(0, 4)$.\\
Reescribimos la recta: $11x - 12y - 5 = 0$. Aquí $A=11, B=-12, C=-5$.
\[
d = \frac{|11(0) - 12(4) - 5|}{\sqrt{(11)^2 + (-12)^2}} = \frac{|-48 - 5|}{\sqrt{121 + 144}} = \frac{|-53|}{\sqrt{265}} = \frac{53}{\sqrt{265}} = \frac{53\sqrt{265}}{265}
\]

\vspace{0.5cm}
\textbf{35.} Encuentre la distancia entre la recta $2x - y = 6$ y el punto de intersección de las rectas $3x - 2y = 1$ y $6x + 3y = 32$.\\
\textit{Paso 1: Encontrar el punto de intersección.}\\
Tenemos el sistema de ecuaciones:
\begin{align*}
    3x - 2y &= 1 \quad \text{--- (Ecuación 1)} \\
    6x + 3y &= 32 \quad \text{--- (Ecuación 2)}
\end{align*}
Multiplicamos la Ecuación 1 por $-2$:
\[
-6x + 4y = -2
\]
Sumamos esto a la Ecuación 2:
\[
(-6x + 4y) + (6x + 3y) = -2 + 32 \implies 7y = 30 \implies y = \frac{30}{7}
\]
Sustituimos $y$ en la Ecuación 1:
\[
3x - 2\left(\frac{30}{7}\right) = 1 \implies 3x - \frac{60}{7} = \frac{7}{7} \implies 3x = \frac{67}{7} \implies x = \frac{67}{21}
\]
El punto de intersección es $P\left(\dfrac{67}{21}, \dfrac{30}{7}\right)$.

\vspace{0.25cm}
\textit{Paso 2: Calcular la distancia de $P$ a la recta $2x - y - 6 = 0$.}\\
Aquí $A=2, B=-1, C=-6$.
\[
d = \frac{\left|2\left(\frac{67}{21}\right) - 1\left(\frac{30}{7}\right) - 6\right|}{\sqrt{(2)^2 + (-1)^2}} = \frac{\left|\frac{134}{21} - \frac{90}{21} - \frac{126}{21}\right|}{\sqrt{4 + 1}} = \frac{\left|\frac{134 - 216}{21}\right|}{\sqrt{5}} = \frac{\left|-\frac{82}{21}\right|}{\sqrt{5}}
\]
\[
d = \frac{82}{21\sqrt{5}} = \frac{82\sqrt{5}}{105}
\]

\end{document}
